
\section{Introduction}
\label{sect:intro}

Every algebra is a quotient of a free algebra. This fundamental result establishes that any algebraic structure, like a monoid, group, or ring, can be presented by generators and relations.
% 
Roughly, this means that, in order to describe an algebra, we just have to pick a set of basic elements, construct the free algebra of formal expressions built upon them, and then impose equations that define the relationships between these expressions. The desired algebra is then obtained by taking the quotient of that free algebra by the specified equations.
% 
Furthermore, we can present homomorphisms of algebras in a similar way. 
That is, given presentations of two algebras, we can describe a homomorphism between them 
as a homomorphism between the associated free algebras which preserves the specified equations. 

This fact finds a general and uniform explanation using the language of category theory, in particular, the notions of an (Barr) exact category \cite{Barr71} and the associated exact completion \cite{CarboniA:freecl, CarboniA:regec}.
% 
Exact categories are categories with finite limits and a well-behaved notion of quotient for internal equivalence relations. Among them are abelian categories and categories of algebras over $\Set$ (the category of sets and functions).
% 
The exact completion is a construction that turns any category with weak finite limits into an exact one. Its objects are pairs of an object and a (pseudo)equivalence relation in the original category, and its arrows are classes of arrows\footnote{Two parallel arrows belong to the same class if, intuitively, they map equivalent elements of their domain to equivalent elements of their codomain.} between the underlying objects that preserve the equivalence relations.
% 
A fundamental result about the exact completion states that the category of algebras for a monad on an exact category where (regular) epimorphisms split  is the exact completion of its subcategory of free algebras \cite{Vitale94}. This result encompasses categories of algebras over $\Set$ and captures, in categorical terms, the fact that algebras can be presented by generators and relations.

The distinguishing feature of exact categories that makes this analysis possible is the ability of computing quotients. Relational doctrines \cite{DagninoP23, DagninoP25tac,DagninoP25apal} have been recently introduced as a framework supporting a very general notion of quotient, together with an associated universal construction dubbed \emph{relational quotient completion}. \footnote{This is inspired by and generalizes the elementary quotient completion of an existential elementary doctrine by Maietti and Rosolini~\cite{MaiettiME:eleqc}.} As specific instances of these notions one recovers exact categories and the exact completion, as well as many other categories of equivalence relations, such as equilogical spaces and categories of assemblies for partial combinatory algebras \cite{lmcs:4302}, but also categories of metric spaces and of normed vector spaces.

In this paper, we study to which extent the aforementioned results about categories of algebras and the exact completion extend to relational doctrines of algebras and the relational quotient completion. This allows us to cover a wider range of algebras and presentations, including, for instance, variants of quantitative algebras~\cite{MardarePP16, MardarePP17, Adamek22}.
% 
To achieve this, we study a notion of projective object for relational doctrines with quotients and characterize the relational doctrines obtained through the relational quotient completion as those admitting a projective cover. 
This generalizes the well-known theorem in \cite{CarboniA:regec} that characterizes the exact categories that are an exact completion as those with enough regular projectives.
%
Then, given that relational doctrines form a 2-category, we can consider monads on relational doctrines defined as monads in such a 2-category \cite{Street72}, which induce a relational doctrine of algebras \cite{DagninoP25tac}. We show that every projective cover of a relational doctrine determines a projective cover of the relational doctrine of algebras for a monad on it, consisting of those free algebras with projective generators. This shows that relational doctrines of algebras arise as relational quotient completions of a doctrine on a suitable subcategory of free algebras. Moreover, when the relational doctrine enjoys a form of the Rule of Choice, analogous to the requirement that epimorphisms split, we have that all free algebras are projective, thus obtaining a result analogous to the one for the exact completion.


The paper is organized as follows.
In \cref{sect:prelim} we recall some basic properties of relational doctrines and the relational quotient completion.
\cref{sect:proj} provides the characterization of those doctrines arising from the relational quotient completion via projective covers,  showing also that the characterisation in \cite{CarboniA:regec} of the exact completion is an instance of our theorem for a specific class of relational doctrines. 
Finally, in \cref{sect:algebra} we study doctrines of algebras for monads on relational doctrines, describing in which cases these arise as relational quotient completions of their restriction to (subcategories of) free algebras, again  extending the analogous result for the exact completion. 
To this end, we also develop some basic results about factorization systems and choice principles in relational doctrines with quotients. 


